%=============================================================
%  第1章 はじめに
%=============================================================
\chapter{はじめに}
\label{chap:intro}

\section{研究背景}
\label{sec:intro-background}

企業の財務構造は静的なものではない。
企業は誕生から成長，成熟，そして衰退に至るライフサイクルの各段階において，
資産構成，負債構造，収益パターン，キャッシュフローの構造を体系的に変化させる。
\tcite{Mueller1972} は産業組織論の観点から企業ライフサイクル理論を提唱し，
経営者の投資行動と利益分配がライフサイクル段階に依存することを理論的に示した。
その後，\tcite{Dickinson2011} は営業・投資・財務キャッシュフローの
符号パターンに基づく5段階分類（導入期・成長期・成熟期・動揺期・衰退期）を提案し，
企業ライフサイクルの実証的な測定手法として広く採用されるに至った。
\tcite{Habib2019survey} のサーベイが示すように，この枠組みは現在，
会計・財務・コーポレートガバナンスの各分野における研究の基盤として機能している。

ライフサイクル理論によれば，企業の財務政策は段階に応じて体系的に変化する。
\tcite{DeAngelo2006} は配当政策がライフサイクル段階に依存することを，
\tcite{Hirsch2011} は負債水準が段階に応じて変化することを示した。
\tcite{Faff2016} は配当・投資・資金調達・現金保有の4つの財務政策が
ライフサイクルに沿って共変動することを包括的に実証している。
こうした知見は，企業の財務構造が単一の指標ではなく
多変量的な特性として捉えるべきであること，
そしてその特性がライフサイクルに沿って動態的に変化することを示唆する。

日本経済は本研究の分析期間（2010〜2024年度）において，
複数の構造的変化を経験した。
1990年代のバブル崩壊後に始まった長期停滞（いわゆる「失われた30年」）の中で，
日本企業は内部留保の蓄積と負債の圧縮を進め，
財務構造は全体として保守化の傾向を示してきた。
2012年末に始まったアベノミクスと日本銀行による異次元金融緩和（量的・質的金融緩和）は，
企業の資金調達環境を大きく変容させた。
超低金利環境の長期化は企業の借入コストを歴史的低水準に押し下げ，
負債構造に影響を与えたと考えられる。
一方で，コーポレートガバナンス・コードの導入（2015年）は，
上場企業に対してROE目標の設定や政策保有株式の縮減を促し，
資本効率性に対する市場の要請を強めた。

さらに，2020年以降の COVID-19 パンデミックは，
業種によって非対称的な影響をもたらした。
飲食・宿泊・運輸業など対面サービス業は急激な収益悪化に直面する一方，
情報通信業やEC関連企業はデジタルトランスフォーメーション（DX）の
加速により成長機会を獲得した。
こうしたマクロ経済環境の変動と制度変更は，企業の財務構造の変化パターン
――すなわち，どのような企業がどのような財務状態に遷移するか――に
影響を与えている可能性がある。
本研究が対象とする15年間は，このように多様な構造変化を含む期間であり，
企業の財務構造ダイナミクスを分析するうえで豊富な変動情報を提供する。

しかし，既存の企業ライフサイクル研究および財務分類研究の多くは，
ある時点における企業の分類に焦点を当てており，
企業が年度をまたいでどのように財務構造を変化させるか
――すなわち，クラスター間の遷移――については十分に解明されていない。
\tcite{GonzalezRuiz2021} は K-Means クラスタリングとマルコフ連鎖を用いた遷移分析を
試みた先駆的研究であるが，対象は 35 社に限定されており，
遷移を引き起こす要因の体系的な解明には至っていない。
また，既存研究の多くは米国や欧州の企業を対象としており，
日本上場企業の財務構造遷移に関する大規模な実証分析は行われていない。

日本の上場企業は，有価証券報告書を通じて継続的かつ標準化された形式で
財務データが開示されている。
金融商品取引法に基づくこの開示制度は，
企業横断的・時系列的に整合性のある財務データを提供しており，
パネルデータ分析の基盤として優れた環境を形成している。
約 \num{3558} 社・15 年間の大規模パネルデータを活用することで，
米国や欧州を対象とした既存研究を補完し，
日本企業に固有の財務構造ダイナミクスを解明する実証分析が可能となる。

近年，機械学習手法の財務分析への応用が急速に進展している。
特に，勾配ブースティング手法である XGBoost \citep{ChenGuestrin2016} は，
倒産予測をはじめとする財務予測タスクにおいて伝統的な統計モデルを上回る
予測精度を達成することが実証されている \citep{Nguyen2023, BenJabeur2023}。
また，SHAP（Shapley Additive Explanations; \citealt{LundbergLee2017}）との
組み合わせにより，高い予測精度を維持しつつモデルの解釈可能性を確保する
フレームワークが確立されつつある \citep{Zhang2023}。
SHAP は協力ゲーム理論における Shapley 値 \citep{Shapley1953} に基礎を置き，
各特徴量の予測への貢献度を理論的に正当化された方法で定量化する。
これらの手法を財務クラスター遷移の分析に応用することで，
どの財務変数が遷移と関連し，遷移先の予測に寄与するかを定量的に解明できる可能性がある。

\section{研究課題}
\label{sec:intro-questions}

以上の背景を踏まえ，本研究では以下の3つの研究課題（RQ）を設定する。

\begin{description}[style=nextline, leftmargin=2em]
  \item[\textbf{RQ1}] 日本上場企業の財務構造は，年度をまたいでどのようなクラスターを形成し，
    各クラスターはどの程度安定的に維持されるか。
\end{description}

企業のライフサイクルに関する理論的予測によれば，
財務構造は段階的に変化するため，ある程度の安定性を示すことが期待される。
しかし，\tcite{Dickinson2011} のキャッシュフロー分類とは異なり，
11の財務変数に基づくデータ駆動型の分類において
どの程度の安定性が観察されるかは実証的な問いである。
クラスターの安定性の程度は，財務構造が戦略的・構造的に決定されているのか，
それとも短期的なショックに応じて頻繁に変化するのかを識別する手がかりとなる。

\begin{description}[style=nextline, leftmargin=2em]
  \item[\textbf{RQ2}] クラスター間の遷移（財務構造の変化）は，
    どの財務変数と関連しているか。
\end{description}

企業がある財務クラスターから別のクラスターへ移行する際に，
どの財務指標の変化がその遷移と関連しているかを明らかにすることは，
ライフサイクル理論の実証的な検証に寄与する。
\tcite{DeAngelo2006} は配当政策が，\tcite{Hirsch2011} は負債水準が，
それぞれライフサイクル段階と関連することを示しているが，
複数の財務変数を同時に考慮した場合にどの変数が相対的に重要であるかは未解明である。

\begin{description}[style=nextline, leftmargin=2em]
  \item[\textbf{RQ3}] 遷移が生じた場合，遷移先クラスターはどの財務変数によって
    予測されるか。また，遷移先によって関連する要因は異なるか。
\end{description}

遷移の有無だけでなく，遷移先を予測できるかどうかは，
経営戦略や財務政策の観点から重要な問いである。
財務脆弱クラスターへの遷移と大企業クラスターへの遷移では，
関連する財務変数が異なる可能性がある。
こうした非対称性の存在は，企業の財務構造変化が単一のメカニズムでは
説明できないことを示唆し，遷移先に応じた差別化された理解の必要性を提起する。

RQ1 は PCA + K-Means クラスタリングとハンガリアン法によるラベル整合
（第\ref{chap:data-method}章）を通じて検討する。
RQ2 は二値プロビット分析および XGBoost の二値分類モデル + SHAP 分析により，
RQ3 は多項ロジット分析および XGBoost の多クラス分類モデル + SHAP 分析により，
それぞれ検討する。

\section{本研究の貢献}
\label{sec:intro-contribution}

本研究は，以下の3点で既存研究に対して貢献する。

\begin{enumerate}
  \item \textbf{方法論的貢献――既存手法の創意的な組み合わせによる
    時系列的追跡可能性の実現}：
    PCA による次元削減，K-Means クラスタリング \citep{Lloyd1982}，
    およびハンガリアン法 \citep{Kuhn1955} による年度間ラベル整合を組み合わせることで，
    年度をまたいで追跡可能な財務クラスターを定義する手法を構築した。
    各手法は確立された技術であるが，これらを統合して年度間の
    クラスター対応関係を客観的に確立する枠組みは，
    財務クラスター遷移分析の文脈では新規の試みである。

  \item \textbf{実証的貢献――日本上場企業の財務クラスター構造と
    遷移パターンの記述}：
    日本上場企業 \num{3558} 社・15 年間のパネルデータを用いて，
    財務クラスター遷移の実態を体系的に記述した。
    分析の結果，5つの財務クラスターが識別され，
    年度間の自己遷移率が高い水準にあることが明らかになった。
    この強いパス依存性は，企業の財務構造が短期的なショックよりも
    構造的な要因によって規定されていることを示唆する。
    なお，二値プロビット分析における説明力（Pseudo $R^2$ = 0.043）は低く，
    個々の財務変数の水準のみからでは遷移の有無を十分に説明できないことも示された。
    この結果は，遷移の発生が財務変数の水準よりも変化の方向やタイミングに
    依存している可能性を示唆する。

  \item \textbf{予測的貢献――XGBoost + SHAP による遷移要因の定量的解明}：
    XGBoost と SHAP を用いた2段階予測モデルにより，
    遷移の有無と遷移先の双方について，関連する財務変数を特定した。
    SHAP 分析により，遷移先クラスターごとに異なる財務変数が予測に寄与すること
    ――遷移の非対称性――が明らかになった。
    この非対称性は，企業の財務構造変化が多元的なメカニズムに
    駆動されていることを示唆する。
    ただし，SHAP 値は予測への貢献度を示すものであり，
    因果関係を意味するものではない点に留意が必要である。
\end{enumerate}

\section{論文の構成}
\label{sec:intro-structure}

本論文の構成は以下のとおりである。

第\ref{chap:literature}章では，企業ライフサイクル理論，財務政策とライフサイクルの関係，
財務クラスタリングと機械学習に関する先行研究を概観し，
本研究の位置づけと既存研究に対するギャップを明らかにする。

第\ref{chap:data-method}章では，分析に用いるデータ（日本上場企業 \num{3558} 社，
2010〜2024年度）の概要と前処理手順，および PCA + K-Means クラスタリング，
ハンガリアン法によるラベル整合，二値プロビット，多項ロジット，
XGBoost + SHAP の各分析手法を詳述する。

第\ref{chap:results}章では，5クラスターの特性，遷移行列，
二値・多クラス予測モデルの性能，SHAP による特徴量重要度，
およびロバストネスチェックの結果を報告する。

第\ref{chap:discussion}章では，各クラスターの経済的解釈，
パス依存性のメカニズム，遷移の非対称性について先行研究と対比しつつ考察し，
投資家・経営者・政策立案者への含意を論じる。

第\ref{chap:conclusion}章では，研究成果を要約し，学術的貢献を整理した上で，
本研究の限界と今後の研究方向を述べる。
